\documentclass[11pt]{article}

\usepackage[T1]{fontenc}
\usepackage{lmodern}
\usepackage{amsmath,amssymb,amsthm,mathtools}
\usepackage[margin=1in]{geometry}
\usepackage{booktabs}
\usepackage{longtable}
\usepackage{array}
\usepackage{microtype}
\usepackage{xcolor}
\usepackage{enumitem}
\usepackage{listings}
\usepackage[hidelinks]{hyperref}

\hypersetup{
  pdftitle={Semiregular C37 Conference Lifts at Order 334},
  pdfauthor={Alec Kriebel, with heavy assistance from ChatGPT 5.6 Sol},
  pdfsubject={Quotient classification, modular realizations, and low-rank obstructions for a semiregular conference-matrix lift},
  pdfkeywords={Hadamard matrix, conference matrix, strongly regular graph, semiregular cyclic action, C37, exhaustive enumeration},
  pdfcreator={LaTeX}
}

\definecolor{codegray}{gray}{0.96}
\lstset{
  basicstyle=\ttfamily\small,
  backgroundcolor=\color{codegray},
  frame=single,
  columns=fullflexible,
  keepspaces=true,
  breaklines=true,
  showstringspaces=false
}

\newtheorem{theorem}{Theorem}[section]
\newtheorem{proposition}[theorem]{Proposition}
\newtheorem{lemma}[theorem]{Lemma}
\newtheorem{corollary}[theorem]{Corollary}
\theoremstyle{definition}
\newtheorem{definition}[theorem]{Definition}
\newtheorem{remark}[theorem]{Remark}

\newcommand{\F}{\mathbb{F}}
\newcommand{\Z}{\mathbb{Z}}
\newcommand{\QR}{\mathrm{QR}}
\newcommand{\NQR}{\mathrm{NQR}}
\newcommand{\rank}{\operatorname{rank}}
\newcommand{\diag}{\operatorname{diag}}
\newcommand{\Tr}{\operatorname{Tr}}
\newcommand{\im}{\operatorname{im}}
\newcommand{\Aut}{\operatorname{Aut}}
\newcommand{\wt}{\operatorname{wt}}
\newcommand{\one}{\mathbf{1}}
\newcommand{\J}{J}
\newcommand{\e}{\mathrm{e}}

\title{\textbf{Semiregular \(C_{37}\) Conference Lifts at Order \(334\):}\\
Quotient Classification, Modular Realizations, and Low-Rank Obstructions}
\author{Alec Kriebel\\[2pt]\large with heavy assistance from ChatGPT 5.6 Sol}
\date{Research draft prepared 25 July 2026\\Not peer reviewed}

\begin{document}
\maketitle

\begin{abstract}
We study symmetric conference matrices of order \(334\) whose
\(333\)-vertex core admits a semiregular cyclic group of order \(37\),
with nine vertex orbits.  Such a matrix would yield a Hadamard matrix of
order \(668\), but no such matrix is constructed here.

The zero-frequency orbit equations are classified exhaustively.  There
are \(196{,}560{,}000\) labeled orbit-sum matrices, \(625\) classes under
simultaneous permutation, and \(314\) classes after also identifying
global sign.  Their reductions modulo two collapse to three parity
classes, or two up to complement.  No quotient has zero diagonal, and a
Fourier trace argument therefore forces a universal \(6/3\) incidence
law on the diagonal circulant blocks.

We count exact finite-field relaxations.  In characteristic two the
nontrivial Fourier component is a rank-four Hermitian projection in a
nine-dimensional unitary space over \(\F_{2^{36}}/\F_{2^{18}}\); the
orientation-fixed relaxation has between \(2^{719}\) and \(2^{720}\)
points.  In characteristic three the corresponding Hermitian
square-zero census lies between \(2^{1141}\) and \(2^{1142}\) per fixed
quotient.  We also construct an exact-margin binary support for one
integral representative of each of the two complement classes.  Both
supports satisfy the complete graph equation modulo two, equivalently
the conference-core equation modulo eight.

In characteristic \(37\), the first-moment system has rank \(16\), and
for a certified quotient every admissible first moment has an explicit
formal completion through all \(37\) truncated layers.  Physical
\(0/1\)-realizability remains the obstruction.  Within the resulting
exponential family, all constant symmetric generators of ranks at most
three are excluded across all \(625\) quotients.  At the first
nonconstant layer we obtain a \(20+20\)-parameter normal form and exclude
both pure rank-two higher terms and the fully enumerated common nondegenerate
two-plane pencil.  Finally, exhaustive four-cycle and small
semiregular-transvection gates fail to repair either frozen
characteristic-two support.

The retained supports do not satisfy the adjacency equation modulo four:
the best recorded carry has \(672\) nonzero coordinates among \(1503\).
Consequently this paper neither constructs nor excludes a conference
graph, and it does not settle the Hadamard problem at order \(668\).
\end{abstract}

\section{Introduction}

A normalized symmetric conference matrix of order \(334\) may be written
\[
 C=
 \begin{pmatrix}
 0&\one^{T}\\
 \one&S
 \end{pmatrix},
 \qquad C^2=333I_{334}.
\]
Its symmetric \(333\times333\) core satisfies
\[
 S\one=0,\qquad S^2=333I_{333}-J_{333}.
 \tag{1.1}\label{eq:seidel-core}
\]
Equivalently,
\[
 A=\frac{J-I-S}{2}
\]
would be a strongly regular graph with parameters
\[
 (v,k,\lambda,\mu)=(333,166,82,83).
 \tag{1.2}\label{eq:srg-parameters}
\]
The standard conference doubling construction would then produce a
Hadamard matrix of order \(668\); the conference-matrix correspondence and
doubling construction are classical \cite{goethals-seidel}.

At the date of this draft, order \(668\) remains listed as the smallest
unresolved Hadamard order in both a recent construction database and a
public open-problem list \cite{cati-pasechnik,epoch-hadamard}, and the
strongly regular graph parameter set \eqref{eq:srg-parameters} remains open
in the standard table \cite{brouwer-srg}.  Mathon's title ``Symmetric
conference matrices of order \(pq^2+1\)'' is a potentially misleading
near-match because \(333=37\cdot3^2\).  In his Theorem~4.1, however, the
same integer \(t\) determines \(q=4t-1\) and \(p=4t+1\), so \(p=q+2\);
the pair \((p,q)=(37,3)\) is outside that construction \cite{mathon}.

The orbit-matrix method for strongly regular graphs with prime-order
automorphisms is classical \cite{behbahani-lam}, as is the general
multicirculant or \(m\)-Cayley framework \cite{martinez}.  A recent
multiplier-obstruction theorem for Legendre pairs of length \(333\)
concerns a different, common-multiplier route and explicitly leaves
unrestricted existence open \cite{ramos-hulak-queiroz}.  Our claims concern
the stated nine-orbit calculation and its quotient-specific consequences,
not those general frameworks.

We impose a semiregular \(C_{37}\)-action with nine vertex orbits.  This is
strictly broader than a regular action by a group of order \(333\).  The
result is a finite but still enormous \(45\)-block lifting problem.  The
main contribution is to identify exactly what is already forced, what can
be realized modulo small primes, and which natural formal and switching
families are impossible.

\subsection*{Main results}

The principal results are as follows.

\begin{enumerate}[label=(\roman*)]
\item The integral zero-frequency quotient system has exactly \(625\)
permutation classes and \(314\) sign-permutation classes
(Theorem~\ref{thm:quotient-census}).
\item The quotient census has only three parity types, two up to
complement, and every integral lift satisfies a \(6/3\)
quadratic-residue incidence law on its diagonal blocks
(Theorems~\ref{thm:parity-types} and \ref{thm:trace-law}).
\item The orientation-fixed characteristic-two relaxation is a
unitary Grassmannian of size strictly between \(2^{719}\) and \(2^{720}\).
The characteristic-three relaxation is counted separately
(Theorems~\ref{thm:mod2-count} and \ref{thm:mod3-count}).
\item Both complement classes have explicit exact-margin binary supports
satisfying every equation modulo two (Theorem~\ref{thm:char2-supports}).
\item The characteristic-\(37\) first-moment map has rank \(16\).  For one
certified quotient every admissible first moment extends through all
truncated layers to a formal solution, but not necessarily to binary blocks
(Theorems~\ref{thm:first-moment} and \ref{thm:formal-completion}).
\item Constant generators of ranks one, two, and three, as well as the
diagonal constant family, are impossible in the explicit formal
conjugation family.  Two fully enumerated rank-two subfamilies at the first
nonconstant layer are also impossible
(Theorems~\ref{thm:constant-low-rank} and
\ref{thm:first-nonconstant-exclusions}).
\item Neither ordinary four-cycle switching nor the smallest
semiregular unitary-transvection family repairs the two frozen modular
supports (Theorem~\ref{thm:switch-gates}).
\end{enumerate}

Computational assertions are tied to deterministic replay programs or
finite certificates, with separate checks for several key counts;
Section~\ref{sec:reproducibility} gives a map from claims to artifacts.

\subsection*{Scope}

The distinction between a relaxation and a construction is essential.
The two binary supports satisfy the adjacency equation only modulo two.
None satisfies it modulo four, and no integer conference graph, symmetric
conference matrix, or Hadamard matrix is produced.  The formal
characteristic-\(37\) completion similarly solves a matrix equation in a
truncated group algebra while deliberately forgetting the physical
\(0/1\) coefficient alphabet.  All results are restricted to the
semiregular \(C_{37}\) lane; they do not exclude that lane, much less
conference graphs without such an automorphism.

\section{The semiregular group-ring model}
\label{sec:model}

Let
\[
 U(x)=1+x+\cdots+x^{36}
\]
in the group ring \(\Z[C_{37}]\cong\Z[x]/(x^{37}-1)\).  A graph invariant
under a semiregular \(C_{37}\)-action on nine fibers is represented by a
\(9\times9\) matrix
\[
 D(x)=(d_{ij}(x))
\]
of binary cyclic words.  Undirectedness and looplessness give
\[
 d_{ji}(x)=d_{ij}(x^{-1}),\qquad [x^0]d_{ii}(x)=0.
 \tag{2.1}\label{eq:star-symmetry}
\]
The strongly regular equation is
\[
 D(x)^2+D(x)=83I_9+83U(x)J_9.
 \tag{2.2}\label{eq:group-ring-srg}
\]

The augmentation
\[
 B=D(1)
\]
is the equitable adjacency quotient.  In Seidel coordinates define
\[
 T=37J_9-I_9-2B.
 \tag{2.3}\label{eq:T-from-B}
\]
Then \(T\) is the matrix of block row sums of the core \(S\).  Evaluation
of \eqref{eq:seidel-core} at the trivial character gives
\[
 T\one=0,\qquad T^2=333I_9-37J_9.
 \tag{2.4}\label{eq:quotient-system}
\]
The diagonal entries of \(T\) are even sums of \(36\) signs; the
off-diagonal entries are odd sums of \(37\) signs.

\begin{lemma}[Local quotient rigidity]
\label{lem:local-rigidity}
Every solution of \eqref{eq:quotient-system} satisfying the block bounds
has
\[
 T_{ii}\in\{-16,-12,-8,-4,0,4,8,12,16\},
 \qquad |T_{ij}|\leq15\quad(i\neq j).
\]
\end{lemma}

\begin{proof}
For one row, write \(d=T_{ii}\) and \(x_1,\ldots,x_8\) for its odd
off-diagonal entries.  The row equations are
\[
 d+\sum_jx_j=0,\qquad d^2+\sum_jx_j^2=296.
\]
Every odd square is \(1\bmod8\), so \(d^2\equiv0\bmod8\), whence
\(4\mid d\).  The norm equation gives \(|d|\leq16\).  If some
\(|x_j|=17\), the remaining seven odd squares and \(d^2\) must contribute
only \(7\), forcing \(d=0\) and the seven remaining entries to be
\(\pm1\); they cannot cancel \(\pm17\).
\end{proof}

There are exactly \(411\) unordered row multisets satisfying these two
row equations, expanding to \(2{,}109{,}524\) ordered rows.  This finite
row catalog is the input to the quotient census.

\section{Complete quotient classification}
\label{sec:quotients}

\begin{theorem}[Exact quotient census]
\label{thm:quotient-census}
The system \eqref{eq:quotient-system}, together with symmetry, parity, and
the block bounds, has exactly
\[
\begin{array}{rcl}
196{,}560{,}000&&\text{fully labeled solutions},\\
625&&\text{classes under }T\sim PTP^T,\\
314&&\text{classes under }T\sim\pm PTP^T.
\end{array}
\]
Exactly three permutation classes are equivalent to their negatives.
The automorphism-order distribution is
\[
\begin{array}{c|rrrrrr}
|\Aut(T)|&1&2&3&4&6&24\\ \hline
\#\text{ classes}&480&100&24&5&14&2.
\end{array}
\]
Every quotient has rank \(4\) over \(\F_{37}\).  There are \(111\)
diagonal multisets, none equal to \(0^9\).
\end{theorem}

\begin{proof}[Certified enumeration]
The enumerator fills a symmetric matrix row by row from the complete
ordered row catalog.  A new row agrees with the previously fixed upper
triangle and has inner product \(-37\) with every earlier row.  The root
row is sorted.  At later depths, unprocessed vertices with identical
columns against earlier rows form interchangeable cells, and new entries
inside each cell are required to be nondecreasing.  This is an orderly
augmentation in the standard isomorph-free generation sense
\cite{mckay}: every solution has a relabeling satisfying the imposed
order.

There are \(7016\) orderly terminal matrices.  Each terminal is checked
directly against every entry of \eqref{eq:quotient-system}, canonicalized
under all \(9!\) permutations, and assigned a stabilizer.  The labeled
total is also reconstructed from the orbit-stabilizer identity
\[
 \sum_{[T]}\frac{9!}{|\Aut(T)|}=196{,}560{,}000.
\]
The packaged certificate
\path{z37_quotient_census_canonical_625.txt} contains the \(625\)
canonical upper triangles together with the audit summaries.  Its
SHA-256 is
\[
\texttt{c5d8765da49deb39c2ff3407b9d0f265e3ca56c1015d5b0075355c53ca60fb5b}.
\]

For the rank assertion, \(T^2=0\bmod37\) makes \(\im T\) totally
isotropic in dimension nine.  Hence \(\rank T\leq4\).  On
\(\one^\perp\), the determinant has \(37\)-adic valuation \(4\).  The
decomposition
\[
\F_{37}^9=\langle\one\rangle\oplus\one^\perp
\]
therefore gives \(\rank T\geq4\).
\end{proof}

An explicit feasible quotient is
\[
\left(
\begin{array}{rrrrrrrrr}
0&-11&-3&7&-3&7&-3&7&-1\\
-11&-4&7&-1&7&-1&7&-1&-3\\
-3&7&0&-11&-3&7&-3&7&-1\\
7&-1&-11&-4&7&-1&7&-1&-3\\
-3&7&-3&7&0&-11&-3&7&-1\\
7&-1&7&-1&-11&-4&7&-1&-3\\
-3&7&-3&7&-3&7&0&-11&-1\\
7&-1&7&-1&7&-1&-11&-4&-3\\
-1&-3&-1&-3&-1&-3&-1&-3&16
\end{array}
\right).
\tag{3.1}\label{eq:certified-T}
\]
It is signed-permutation equivalent to a \(1+4+4\) normal form and is one
of the two order-\(24\) permutation classes.  Those two classes are
negatives, hence form the unique maximally symmetric sign-class.

\begin{theorem}[Parity collapse]
\label{thm:parity-types}
Reducing \(B=(37J-I-T)/2\) modulo two maps the \(625\) integral quotient
classes onto exactly three binary graph classes.  Their integral preimage
counts are
\[
206,\quad213,\quad206,
\]
and their automorphism orders are
\[
24,\quad20,\quad24.
\]
The first and third are complements; the middle class is
self-complementary up to permutation.  Thus two types remain after
identifying complement.
\end{theorem}

The number of fully labeled binary quotients is consequently
\[
\frac{9!}{24}+\frac{9!}{20}+\frac{9!}{24}=48{,}384.
\]
One order-\(24\) representative consists of a \(K_4\), an independent
set of order four, a ninth vertex joined to the independent set, and a
matching between the two four-sets.  Its degree sequence is
\(2^4,4^5\).  Its complement has degree sequence \(4^5,6^4\), while the
order-\(20\) class has degree sequence \(2,2,4,4,4,4,4,6,6\).

\section{The universal diagonal trace law}
\label{sec:trace-law}

For \(t\in\F_{37}^{\times}\), let \(g(t)\) be the number of the nine
diagonal adjacency blocks containing displacement \(t\).

\begin{theorem}[The \(6/3\) law]
\label{thm:trace-law}
After possibly interchanging quadratic residues and nonresidues, every
semiregular \(C_{37}\) lift of every quotient in
Theorem~\ref{thm:quotient-census} satisfies
\[
 g(t)=
 \begin{cases}
 6,&t\in\QR,\\
 3,&t\in\NQR.
 \end{cases}
\tag{4.1}\label{eq:63-law}
\]
\end{theorem}

\begin{proof}
At a nontrivial character \(\zeta^a\), the Hermitian Fourier block
\(D(\zeta^a)\) satisfies
\[
 D(\zeta^a)^2+D(\zeta^a)=83I_9.
\]
Its only eigenvalues are
\[
 r,s=\frac{-1\pm3\sqrt{37}}2.
\]
At frequency zero, \(B=D(1)\) has eigenvalue \(166\) on
\(\langle\one\rangle\).  On \(\one^\perp\), equation
\eqref{eq:quotient-system} gives the conjugate eigenvalues \(r,s\), each
with multiplicity four.  The full
\(\operatorname{srg}(333,166,82,83)\) spectrum has \(r\) and \(s\), each
with multiplicity \(166\).

Cyclotomic Galois conjugacy makes the multiplicity of \(r\) constant,
say \(m\), on the \(18\) quadratic-residue frequencies and constant, say
\(m'\), on the \(18\) nonresidue frequencies.  Counting the full
\(r\)-eigenspace, including the four copies at frequency zero, gives
\[
 4+18m+18m'=166,
\]
so \(m+m'=9\), or \(m'=9-m\).  Consequently the two nonzero-frequency
traces are
\[
 \tau_{\QR}=-\frac92+\frac{3\sqrt{37}}2(2m-9),
 \qquad
 \tau_{\NQR}=-\frac92-\frac{3\sqrt{37}}2(2m-9),
\]
while \(\tau_0=\Tr B=162\).

For \(t\neq0\), the quadratic Gaussian sums, with a fixed choice of
\(\sqrt{37}\), are
\[
 \sum_{a\in\QR}\zeta^{-at}
   =\frac{-1+\chi(t)\sqrt{37}}2,
 \qquad
 \sum_{a\in\NQR}\zeta^{-at}
   =\frac{-1-\chi(t)\sqrt{37}}2.
\]
Fourier inversion of
\(\Tr D(\zeta^a)=\sum_tg(t)\zeta^{at}\) now gives
\[
 g(t)=\frac1{37}
 \left(
 162+\tau_{\QR}\sum_{a\in\QR}\zeta^{-at}
     +\tau_{\NQR}\sum_{a\in\NQR}\zeta^{-at}
 \right)
 =\frac{9+3(2m-9)\chi(t)}2.
\tag{4.2}\label{eq:g-formula}
\]
The pointwise bounds \(0\leq g(t)\leq9\) leave
\(m\in\{3,4,5,6\}\), and therefore, up to orientation, incidence pairs
\((0,9)\) or \((3,6)\).

In the \((0,9)\) branch every diagonal block is one complete quadratic
class and has degree \(18\).  Equation \eqref{eq:T-from-B} would then give
\(T_{ii}=0\) for all \(i\), contradicting the exhaustive assertion in
Theorem~\ref{thm:quotient-census}.  Hence only \((3,6)\) remains.
\end{proof}

\subsection{Ambient search sizes}

There are \(1494\) raw block-membership bits: \(1332\) from the \(36\)
unordered off-diagonal blocks and \(162\) inverse-pair bits from the nine
diagonal blocks.  For one quotient \(T\), the fixed-margin count before
\eqref{eq:63-law} is
\[
 \prod_i
 \binom{18}{(36-T_{ii})/4}
 \prod_{i<j}
 \binom{37}{(37-T_{ij})/2}.
\tag{4.3}\label{eq:margin-count}
\]
Across the \(625\) classes, the extremal base-two logarithms are
\[
1340.52392808143
\quad\text{and}\quad
1340.83193919044.
\]

After the \(6/3\) law and a nonresidue decimation fixing its orientation,
the per-class logarithm lies between
\[
1297.60492221007667
\quad\text{and}\quad
1297.90621474626255.
\]
Summing one canonical space for each quotient gives
\[
\log_2(\text{\(625\)-class trace-law union})
=1307.10873431446430.
\tag{4.4}\label{eq:trace-union}
\]
These are exact combinatorial counts before the remaining Fourier
equations, not estimates of conference graphs.

\section{Exact finite-field relaxations}
\label{sec:finite-field}

\subsection{Characteristic two}

Because \(2\) is primitive modulo \(37\),
\[
 \F_2[C_{37}]\cong \F_2\times K,
 \qquad K=\F_{2^{36}}.
\]
Group inversion on \(K\) is
\[
 \bar a=a^{2^{18}},
\]
with fixed field \(L=\F_{q}\), \(q=2^{18}\).  On the nontrivial factor,
\eqref{eq:group-ring-srg} reduces to
\[
 D^2+D=I.
\tag{5.1}\label{eq:char2-D}
\]
For a primitive \(\omega\in\F_4\subseteq L\), one of the two shifts
\[
 E=D+\omega^2I
\]
is Hermitian and idempotent.  Hence \(E\) is the orthogonal projection
onto a nondegenerate subspace of the unitary space \(K^9\).

Recall
\[
 |U(n,q)|=q^{n(n-1)/2}
 \prod_{i=1}^{n}\bigl(q^i-(-1)^i\bigr).
\tag{5.2}\label{eq:unitary-order}
\]
The finite unitary order formula, Witt geometry, and the associated
transitivity and isotropic-subspace counts used below are standard; see,
for example, Taylor's treatment of finite classical groups
\cite{taylor-classical-groups}.

\begin{theorem}[Characteristic-two census]
\label{thm:mod2-count}
For a fixed parity quotient and a fixed \(6/3\) orientation, the
nontrivial characteristic-two relaxation consists of exactly
\[
 N_4=
 \frac{|U(9,q)|}{|U(4,q)|\,|U(5,q)|},
 \qquad q=2^{18},
\tag{5.3}\label{eq:N4}
\]
rank-four Hermitian projections.  This \(217\)-digit integer satisfies
\[
2^{719}<N_4<2^{720}.
\tag{5.4}\label{eq:N4-bound}
\]
The opposite orientation has rank five and the same count.  Up to fiber
permutation and complement, every integral lift in the complete quotient
census injects into a union of at most two such sets; consequently
\[
2^{720}<2N_4<2^{721}.
\tag{5.5}\label{eq:all-mod2-bound}
\]
\end{theorem}

\begin{proof}
Hermitian idempotents are in bijection with nondegenerate subspaces via
their images.  The unitary group is transitive on nondegenerate
\(r\)-spaces, with stabilizer \(U(r,q)\times U(9-r,q)\), giving
\eqref{eq:N4}.  The trace law fixes \(r=4\) in one orientation and \(r=5\)
in the other.  Fixing both the trivial CRT factor and the nontrivial field
factor uniquely fixes each binary word of length \(37\), so every integral
lift injects into this relaxation.  The parity collapse in
Theorem~\ref{thm:parity-types} gives the factor two.
\end{proof}

The same calculation gives the characteristic-polynomial identity
\[
\det(\lambda I+D(x))
=\lambda^9+C(x)(\lambda^8+\lambda^4+1)+\lambda^5+\lambda
\pmod2,
\tag{5.6}\label{eq:char2-characteristic}
\]
where \(C\) is the oriented quadratic-class indicator.  Equivalently,
write
\[
\det(\lambda I+D)=\sum_{i=0}^{9}e_i(D)\lambda^{9-i},
\qquad e_0(D)=1,
\]
so \(e_i(D)\in\F_2[C_{37}]\) is the \(i\)th elementary symmetric
coefficient, and let \(\delta=x^0\) be the identity element of this group
algebra.  Then
\[
e_1=e_5=e_9=C,\qquad
e_4=e_8=\delta,\qquad
e_2=e_3=e_6=e_7=0.
\tag{5.7}\label{eq:char2-minors}
\]

\subsection{Characteristic three}

Here
\[
 \F_3[C_{37}]
 \cong
 \F_3\times \F_{3^{18}}\times\F_{3^{18}},
\]
and inversion on either nontrivial factor is the unitary involution over
\(\F_q\), \(q=3^9\).  If \(N=D-I\), the nontrivial equation becomes
\[
 N^2=0,\qquad N^*=N.
\tag{5.8}\label{eq:char3-squarezero}
\]
Its image is totally isotropic, so \(\rank N\leq4\).

Let \(I_{n,r}(q)\) denote the number of totally isotropic \(r\)-spaces in
a nondegenerate \(n\)-dimensional Hermitian space:
\[
I_{n,r}(q)=
\prod_{j=0}^{r-1}
\frac{
\bigl(q^{n-2j}-(-1)^{n-2j}\bigr)
\bigl(q^{n-2j-1}-(-1)^{n-2j-1}\bigr)}
{q^{2(j+1)}-1}.
\tag{5.9}\label{eq:isotropic-count}
\]
Let
\[
H_r(q)=\frac{|GL(r,q^2)|}{|U(r,q)|}
\tag{5.10}\label{eq:hermitian-form-count}
\]
be the number of nonsingular Hermitian forms on a fixed \(r\)-space.

\begin{theorem}[Characteristic-three census]
\label{thm:mod3-count}
For one degree-\(18\) factor, the exact number of Hermitian square-zero
matrices is
\[
 M_3=\sum_{r=0}^{4}I_{9,r}(3^9)H_r(3^9).
\tag{5.11}\label{eq:M3}
\]
For a fixed trivial quotient, the two nontrivial factors are independent,
so the relaxation contains \(M_3^2\) points, with
\[
2^{1141}<M_3^2<2^{1142}.
\tag{5.12}\label{eq:M3-bound}
\]
\end{theorem}

\begin{proof}
A Hermitian square-zero map is determined by its totally isotropic image
\(U\) and the nonsingular Hermitian form induced on \(U\) by the map from
the quotient dual to \(U\).  This gives the factor
\(I_{9,r}(q)H_r(q)\).  The CRT split gives the square.
\end{proof}

\begin{remark}
Neither \eqref{eq:N4-bound} nor \eqref{eq:M3-bound} is an exact count of
binary or integral conference graphs.  These are separate upper bounds
from different modular relaxations and must not be multiplied as if they
were independent filters.
\end{remark}

\section{Exact characteristic-two physical supports}
\label{sec:char2-supports}

The unitary projection census does not itself impose the prescribed
integer Hamming weights of all \(45\) cyclic blocks.  Nevertheless both
complement classes have an explicit exact-margin realization.

\begin{theorem}[Two exact modular supports]
\label{thm:char2-supports}
For one integral quotient representing each of the two parity types up to
complement, there exists a \(9\times9\) matrix of binary cyclic words
satisfying:
\begin{enumerate}[label=(\alph*)]
\item the prescribed exact integer margin of every block;
\item star symmetry and zero diagonal identity coefficients;
\item the \(6/3\) law \eqref{eq:63-law};
\item every one of the \(2997\) ordered group-ring coefficients of
\[
 D(x)^2+D(x)=I_9+U(x)J_9
 \quad\text{over }\F_2.
\tag{6.1}\label{eq:physical-mod2}
\]
\end{enumerate}
Equivalently, the two expanded \(333\times333\) supports are loopless,
have degree \(166\), and satisfy all \(110{,}889\) scalar equations.
\end{theorem}

The text-witness SHA-256 values are
\[
\begin{array}{c|l}
\text{type 1}&
\texttt{3f77fc3d39fc6b8dfd33efbba846e61ca835581d15f9ae4c12d8f57a3249e697}\\
\text{type 2}&
\texttt{2e5087308dceb18398daaba4b1ca5868d755cecb41ae7892f988ae8083c03c22}.
\end{array}
\tag{6.2}\label{eq:support-hashes}
\]

\subsection{Construction outline}

The CRT is particularly explicit.  If \(d\in K\) is represented by a
polynomial of degree below \(36\) and \(q_0\in\F_2\) is the augmentation,
the unique binary cyclic word is
\[
 f=d+c\Phi_{37},\qquad c=q_0+d(1).
\tag{6.3}\label{eq:char2-crt}
\]
A max-flow first prescribes a \(9\times18\) incidence table for diagonal
inverse pairs, simultaneously enforcing all nine margins and column sums
six on residues and three on nonresidues.

Starting from a dense central projection, exact two-coordinate unitary
rotations prescribe eight diagonal entries; the ninth is forced by trace.
The off-diagonal margins are then imposed by diagonal unitary conjugation.
The relevant unit circle has order
\[
2^{18}+1=262{,}145.
\]
For each edge, all phase differences are enumerated, and the remaining
nine-vertex difference constraint is solved exactly.

The following identity is the compatibility gate for off-diagonal norms.

\begin{proposition}[Norm-support trace identity]
\label{prop:norm-support}
Let \(f\) be a binary word of length \(37\) and weight \(k\), evaluated at
a nontrivial \(37\)th root \(\zeta\).  Then
\[
\Tr_{\F_{2^{18}}/\F_2}
\bigl(f(\zeta)f(\zeta^{-1})\bigr)
=\binom{k}{2}\pmod2.
\tag{6.4}\label{eq:norm-support}
\]
\end{proposition}

\begin{proof}
Expand the norm into parity autocorrelations.  The \(18\) inverse-pair
terms have absolute trace one because \(2\) is primitive modulo \(37\).
Their sum counts each unordered pair of support positions once.
\end{proof}

\begin{remark}
Theorem~\ref{thm:char2-supports} concerns one integral representative of
each parity type.  It does not assert exact-margin realizability for all
\(625\) integral quotients.
\end{remark}

\section{The characteristic-\texorpdfstring{\(37\)}{37} moment filtration}
\label{sec:char37}

Work in
\[
 R=\F_{37}[y]/(y^{37}),\qquad x=1+y.
\]
Since \(U(x)=y^{36}\) and \(18^2+18=9\bmod37\), set
\[
 N(y)=D(1+y)-18I_9.
\]
Equation \eqref{eq:group-ring-srg} becomes
\[
 N(y)^2=9y^{36}J_9.
\tag{7.1}\label{eq:yadic-equation}
\]
Its constant term is
\[
 N_0=B-18I_9=-\frac{T}{2}\pmod{37},
\tag{7.2}\label{eq:N0}
\]
a symmetric rank-four square-zero matrix annihilating \(\one\).

\subsection{The first moment}

Star symmetry makes the first coefficient \(N_1\) skew-symmetric.  The
coefficient of \(y\) in \eqref{eq:yadic-equation} is
\[
 N_0N_1+N_1N_0=0.
\tag{7.3}\label{eq:first-moment-system}
\]

\begin{theorem}[Rank-\(16\) first-moment law]
\label{thm:first-moment}
For every quotient in Theorem~\ref{thm:quotient-census}, the linear system
\eqref{eq:first-moment-system} on the \(36\) skew coordinates has rank
\(16\), leaving a \(20\)-dimensional kernel.

For every fixed nontrivial block size, translation acts transitively on
the possible first moments of a subset of \(\F_{37}\).  Hence the
\(16\) equations divide every fixed-margin ambient count by exactly
\[
37^{16}.
\tag{7.4}\label{eq:37^16}
\]
The \(625\)-class trace-law union therefore has logarithm
\[
1307.10873431446430-16\log_2(37)
=1223.75748046440094.
\tag{7.5}\label{eq:first-moment-union}
\]
\end{theorem}

The first three nonconstant binomial moments can be realized
simultaneously with the quotient margins, star symmetry, and the trace
law for the certified quotient \eqref{eq:certified-T}.  The corresponding
physical support satisfies \eqref{eq:yadic-equation} through degree three.
Its first failure is degree four, with \(79\) of the \(81\) matrix entries
nonzero.  Thus several low layers are feasible but are not evidence of
convergence to a graph.

\subsection{A full formal completion}

The correct coordinate for the involution is
\[
 z=\log(1+y),
\]
because cyclic reversal sends \(z\) to \(-z\).  In the truncated ring,
\[
 z^{36}=y^{36}
\quad\text{and}\quad
 z^{18}=6\sum_{t\neq0}\chi(t)x^t,
\tag{7.6}\label{eq:z18-paley}
\]
where \(\chi\) is the quadratic character of \(\F_{37}\).

\begin{theorem}[Formal integrability for the certified quotient]
\label{thm:formal-completion}
Let \(N_0\) be \eqref{eq:N0} for the certified quotient
\eqref{eq:certified-T}.  Every admissible first moment
\[
 X^T=-X,\qquad N_0X+XN_0=0
\]
has the form
\[
 X=[N_0,A]
\]
for some symmetric \(A\in M_9(\F_{37})\).  For either
\(\eta\in\{1,-1\}\),
\[
 N_A(y)=
 \e^{-zA}
 \left(
 N_0+\eta z^{18}J_9+19y^{36}J_9
 \right)
 \e^{zA}
\tag{7.7}\label{eq:formal-family}
\]
has star symmetry and satisfies the full equation
\[
 N_A(y)^2=9y^{36}J_9.
\]
\end{theorem}

\begin{proof}
Exact elimination gives rank \(20\) for the map
\[
\operatorname{Sym}_9(\F_{37})\longrightarrow
\operatorname{Skew}_9(\F_{37}),
\qquad A\longmapsto[N_0,A],
\]
equal to the dimension of the admissible first-moment kernel.
Furthermore,
\[
N_0^2=0,\qquad N_0J=JN_0=0,\qquad J^2=9J.
\]
The term \(19y^{36}J\) lies in the terminal socle: it changes the trace
but neither the square nor lower coefficients.  Since
\((\eta z^{18}J)^2=z^{36}J^2=9y^{36}J\), the middle matrix in
\eqref{eq:formal-family} has the required square.  Conjugation preserves
the square and star symmetry.
\end{proof}

\begin{remark}
Theorem~\ref{thm:formal-completion} is not a classification of all formal
solutions.  More importantly, the coefficients of
\eqref{eq:formal-family} need not be realizable simultaneously by binary
cyclic blocks with the required margins.  The theorem identifies physical
support realizability, rather than formal matrix algebra, as the remaining
obstruction in this chart.
\end{remark}

\section{Constant-generator obstructions through rank three}
\label{sec:constant-rank}

For each diagonal block, let \(r_i(t)\) denote the coefficient after
removing the common terminal \(19y^{36}J\) contribution from
\eqref{eq:formal-family}.  Since \(y^{36}=U=\sum_t x^t\) and
\(D=N+18I\), restoring the physical adjacency coefficient gives
\[
 D_{ii}(0)=r_i(0),\qquad
 D_{ii}(t)=r_i(t)+19\pmod{37}\quad(t\neq0).
\]
Thus the reduced target \(\{18,19\}\) below is equivalent to the restored
adjacency target \(\{0,1\}\).  A generic reduced cyclic word \(r\) must
satisfy
\[
r(0)=0,\qquad r(t)\in\{18,19\}\quad(t\neq0).
\tag{8.1}\label{eq:binary-diagonal-target}
\]
Its number of nonzero-lag \(19\)'s is the diagonal block degree.

\begin{theorem}[Constant low-rank obstruction]
\label{thm:constant-low-rank}
Within the family \eqref{eq:formal-family}:
\begin{enumerate}[label=(\alph*)]
\item no diagonal constant generator \(A\) is physical;
\item no nonzero symmetric rank-one constant generator is physical;
\item no symmetric rank-two constant generator is physical for any of the
\(625\) quotient classes;
\item no symmetric rank-three constant generator is physical for any of
the \(625\) quotient classes.
\end{enumerate}
Thus a viable single constant generator in this family must have rank at
least four.
\end{theorem}

\subsection{Ranks zero, one, and two}

For diagonal \(A\), every restored nonzero-lag coefficient of \(D\) is
\(13\) or \(25\), rather than \(0\) or \(1\); these are not the reduced
\(\{18,19\}\) values in \eqref{eq:binary-diagonal-target}.

Rank one separates into nonisotropic and isotropic cases.  In the
nonisotropic case, six explicitly occurring quadratic-character patterns
force the normalized local parameters to satisfy \(b_i=0\) and
\(a_i\in\{3,-3\}\), while the rank-one parametrization gives
\(b_i=0\Rightarrow a_i=1\).  This contradiction is independent of the
quotient.

In the isotropic case write \(A=\lambda uu^T\) with \(u^Tu=0\), so
\(A^2=0\).  After the terminal \(19y^{36}J\) term has been removed, every
zero-lag diagonal correction is a multiple of \(z^k\) with
\(1\leq k\leq35\).  Exact conversion from the \(z\)-basis to the cyclic
\(x\)-basis gives
\[
 [x^0]z^k=0\qquad(1\leq k\leq35).
\]
Consequently the reduced zero-lag coefficient is fixed, for every
quotient and every coordinate, at
\[
 r_i(0)=N_{0,ii}=B_{ii}-18=-\frac{T_{ii}}2\pmod{37}.
\]
Physicality would require \(r_i(0)=0\) for all \(i\), hence
\(T_{ii}=0\) for all \(i\).  The \(625\)-class canonical census contains
no quotient with all-zero diagonal, so no isotropic rank-one generator is
physical for any quotient.

For rank two, split semisimple types are eliminated by an exhaustive local
character test over all \(35\) projective eigenvalue ratios and all
\(37^2\) projector-coordinate triples.  The remaining rational types are:
\[
\begin{array}{ll}
\text{(i)}&\text{an irreducible quadratic block},\\
\text{(ii)}&\text{a nonzero \(J_2\) block},\\
\text{(iii)}&\text{a nonzero eigenline plus \(J_2(0)\)},\\
\text{(iv)}&J_3(0),\\
\text{(v)}&J_2(0)\oplus J_2(0).
\end{array}
\]
A similarity-invariant temporal overcode gives only the two Paley
weight-\(18\) words in cases (i), (iv), and (v).  Cases (ii) and (iii)
initially allow weights \(16,\ldots,20\), leaving one quotient sign-class
with diagonal Seidel profile
\[
(-4,-4,0,0,0,0,0,4,4).
\]
Restoring the fixed coefficient of \(z^{18}J\) reduces the final diagonal
condition to two affine cosets of a seven-dimensional cyclic code.  The
two genuine orientations \(\eta=\pm1\) have empty binary intersection,
closing the exception.

\subsection{Rank three}

The rank-three audit uses a safe over-list of all projective rational
similarity types, including types not realizable by symmetric matrices:
\[
\begin{array}{l|r}
\text{primary shape}&\text{projective types}\\ \hline
\text{invertible dimension three, cyclic}&1371\\
\text{invertible dimension three, noncyclic}&37\\
J_2(0)\text{ plus invertible dimension two}&39\\
\text{zero-primary rank two plus a nonzero line}&2\\
\text{pure nilpotent rank three}&3\\ \hline
\text{total}&1452.
\end{array}
\tag{8.2}\label{eq:rank3-types}
\]
For a minimal polynomial of degree \(m\leq4\), write
\[
\e^{\pm zA}=\sum_{r=0}^{m-1}f_r^\pm(z)A^r.
\]
Symmetry gives
\[
\diag(A^rMA^s)=\diag(A^sMA^r)
\]
for symmetric \(A,M\).  This yields a reduced temporal code
\[
W_A=F_A+z^{18}F_A
\]
spanned by the symmetrized products of the \(f_r^\pm\).  Its intersection
with \eqref{eq:binary-diagonal-target} contains at most \(256\) words for
any type.  Quotient diagonal profiles remove \(492\) types.  For each of
the remaining \(960\), the audit restores the fixed \(J\) coefficient and
exhausts both trace orientations and at most \(37^3=50{,}653\) relaxed
local power tuples.  No binary word remains.  Every relaxation enlarges
the candidate set, so emptiness proves the stated obstruction.

\begin{remark}
Theorem~\ref{thm:constant-low-rank} does not exclude rank at least four,
a genuinely nonconstant conjugator, a different formal parametrization,
or the unrestricted \(45\)-block lift.
\end{remark}

\section{The first nonconstant conjugator layer}
\label{sec:first-nonconstant}

For a matrix polynomial \(M\), put
\[
M^\dagger=M(-z)^T.
\]
A near-identity conjugator \(G\) preserves star symmetry precisely when
\(G^\dagger=G^{-1}\).  Since logarithm and exponential are inverse on the
nilpotent ideal, write \(G=\e^K\) with \(K^\dagger=-K\).  Thus
\[
K=zA_0+z^2A_1+z^3A_2+\cdots,
\qquad A_r^T=(-1)^rA_r.
\tag{9.1}\label{eq:adjoint-parity}
\]
The first extension beyond a constant generator is
\[
K=zA_0+z^2A_1,
\qquad A_0^T=A_0,\quad A_1^T=-A_1.
\tag{9.2}\label{eq:first-nonconstant-K}
\]

\begin{proposition}[\(20+20\) normal form]
\label{prop:20plus20}
Through degree two, gauge reduction of
\eqref{eq:first-nonconstant-K} leaves \(20\) essential parameters in
\(A_0\) and \(20\) in \(A_1\).
\end{proposition}

\begin{proof}
Expanding
\[
\e^{-K}N_0\e^K=N_0+zX+z^2Y+O(z^3)
\]
gives
\[
X=[N_0,A_0],\qquad
Y=[N_0,A_1]+\frac12[[N_0,A_0],A_0].
\]
The exact commutator ranks are
\[
\begin{array}{rcl}
\operatorname{Sym}_9&\xrightarrow{[N_0,-]}&
\operatorname{Skew}_9:
\quad\rank=20,\quad\dim\ker=25,\\
\operatorname{Skew}_9&\xrightarrow{[N_0,-]}&
\operatorname{Sym}_9:
\quad\rank=20,\quad\dim\ker=16.
\end{array}
\]
The homogeneous second tangent has dimension \(21\); its trace-zero
subspace is exactly the second commutator image.  Centralizer-valued
left multiplication removes the two kernels, leaving the claimed slice.
\end{proof}

\begin{theorem}[Rank-two closures at the first nonconstant layer]
\label{thm:first-nonconstant-exclusions}
Across all \(625\) quotient classes, neither of the following families can
produce binary diagonal blocks:
\begin{enumerate}[label=(\alph*)]
\item the pure first-higher rank-two family \(K=z^2B\), with
\(B^T=-B\) and \(\rank B=2\);
\item the complete family in which \(A_0\) and \(A_1\) in
\eqref{eq:first-nonconstant-K} are supported on one common nondegenerate
two-plane and \(A_1\neq0\).
\end{enumerate}
\end{theorem}

\subsection{Pure first-higher rank two}

Write \(B=uv^T-vu^T\).  Then
\[
B^3=-\Delta B,\qquad
\Delta=(u^Tu)(v^Tv)-(u^Tv)^2,
\]
so there are four projective rational types: split semisimple,
irreducible trace-zero semisimple, \(J_3(0)\), and
\(J_2(0)\oplus J_2(0)\).  In every type, the exact diagonal temporal
overcode meets the binary target only in the two Paley words, both of
weight \(18\).  Theorem~\ref{thm:quotient-census} excludes an all-degree-18
quotient.

\subsection{The common nondegenerate two-plane pencil}

On the active plane, write
\[
A_0=tI+S,\qquad A_1=B,
\]
where \(S^T=S\), \(\Tr S=0\), and \(B^T=-B\neq0\).  Then
\[
S^2=\alpha I,\qquad B^2=\beta I,\qquad SB+BS=0,
\quad\beta\neq0.
\tag{9.3}\label{eq:two-plane-relations}
\]
The traceless exponent closes because
\[
(zS+z^2B)^2=(\alpha z^2+\beta z^4)I.
\tag{9.4}\label{eq:two-plane-square}
\]
After decimation there are \(76\) orbits with \(t=0\) and \(1332\)
normalized parameter types with \(t\neq0\).  A symmetry-restored temporal
overcode leaves only two exceptional nonzero-\(t\) types:
\[
(\alpha,\beta)=(5,32),\quad
(\alpha,\beta)=(19,20),
\]
with allowed weight sets
\[
\{14,18,18,22\},
\qquad
\{12,18,18,24\},
\]
respectively.  Only the second meets any quotient profile, leaving four
quotient classes and \(40\) trace-compatible overcode assignments.

For \((t,\alpha,\beta)=(1,19,20)\), the ordinary function space \(F\)
and \(z^{18}F\) are disjoint seven-dimensional spaces.  Restoring the
fixed \(J=\one\one^T\) coefficient forces, at each coordinate, a
ten-component vector determined by three scalars
\((p,s,b)\in\F_{37}^3\).  Exhausting all \(37^3\) relaxed triples for
both \(\eta=\pm1\) gives zero survivors for every exceptional word.  This
closes the entire named pencil.

\begin{remark}
Degenerate common support with nonzero \(A_0\), support that changes with
degree, and nonconstant coefficients of rank at least three remain open.
\end{remark}

\section{The next \texorpdfstring{\(2\)}{2}-adic digit and small-switch gates}
\label{sec:switches}

Throughout this section \(D\) denotes the expanded \(333\times333\)
binary adjacency matrix represented earlier by \(D(x)\).  Put
\[
 J=J_{333},\qquad I=I_{333},\qquad S=J-I-2D.
\]
Define the expanded carry
\[
R(D)=
\frac{D^2+D-83(I+J)}{2}\pmod2.
\tag{10.1}\label{eq:carry}
\]
Vanishing of \(R(D)\) is exactly the adjacency equation modulo four.
Since
\[
S^2-(333I-J)
=4\bigl(D^2+D-83(I+J)\bigr),
\]
this is equivalently the conference-core equation modulo \(16\).

The two initial supports from Theorem~\ref{thm:char2-supports} have,
among the \(1503\) independent star-symmetric orbit coefficients of this
expanded carry,
\[
\begin{array}{c|cc}
&\text{nonzero carries}&\text{zero carries}\\ \hline
\text{type 1}&722&781\\
\text{type 2}&764&739.
\end{array}
\tag{10.2}\label{eq:initial-carries}
\]
A deterministic \(200{,}000\)-move walk within one exact-margin type-1
phase fiber improves \(722\) to \(672\), consisting of \(72\) diagonal
and \(600\) off-diagonal defects.  This is an attained upper bound, not a
minimum.

\begin{theorem}[Exact small-switch gates]
\label{thm:switch-gates}
For each of the two original frozen supports:
\begin{enumerate}[label=(\alph*)]
\item no ordinary four-cycle toggle preserves the complete
characteristic-two equation;
\item among the \(49{,}284\) tested smallest semiregular Hermitian
transvections, none preserves all exact block margins.
\end{enumerate}
\end{theorem}

\begin{proof}
For a four-cycle toggle,
\[
\Delta=uv^T+vu^T
\]
with disjoint weight-two vectors \(u,v\), and \(\Delta^2=0\).
The toggled matrix preserves \(D^2+D=I+J\) if and only if
\(\langle u,v\rangle\) is \(D\)-invariant.  Exhausting all
\(\binom{333}{2}=55{,}278\) pair-vectors gives minimum
\(\wt(Du)=136\) and \(138\), respectively, and no invariant plane.

For the second family, work in
\(K=\F_2[x]/(\Phi_{37}(x))\) with \(\bar x=x^{-1}\), and take
\[
u=e_a+x^se_b,\qquad a<b,\quad s\in\F_{37},
\]
so \(u^*u=0\).  The audited scalar family is explicitly
\[
\mathcal C_{\mathrm{small}}
=\{1\}\cup
\{x^t+x^{-t}:1\leq t\leq18\}\cup
\{1+x^t+x^{-t}:1\leq t\leq18\}.
\]
Every \(c\in\mathcal C_{\mathrm{small}}\) satisfies \(c=\bar c\), and
\[
 U_c=I_9+c\,uu^*
\]
is unitary.  The audit exhausts the
\(36\cdot37=1332\) binomial isotropic directions and all \(37\) displayed
scalars, hence \(49{,}284\) exact unitary conjugations per support.  All
remain loopless, and
\(1492\) or \(1604\) retain the trace law, but none retains every block
margin.
\end{proof}

These are exclusions of two finite named families, not of all support
switches.  A complete fixed-quotient CP-SAT model for the modulo-four
problem has \(1494\) physical bits, \(418{,}293\) reused Boolean products,
and \(1503\) independent cyclic congruences.  Bounded runs returned
\texttt{UNKNOWN}, including at the lower layer where a witness is known.
Those solver outcomes are diagnostics of the formulation, not
mathematical evidence.

\section{Reproducibility}
\label{sec:reproducibility}

The mathematical source artifacts referenced below are contained in the
repository directory
\begin{center}
\path{hadamard_668_search/conference_334_z37_lift/}.
\end{center}
This paper package also includes the canonical census dump, its digest
sidecar, and a dependency-free structural checker.  The calculations use
exact integer or finite-field arithmetic.  No floating-point comparison,
stochastic nonexistence inference, or timeout is used in a theorem.

\small
\begin{longtable}{>{\raggedright\arraybackslash}p{0.27\textwidth}
                  >{\raggedright\arraybackslash}p{0.66\textwidth}}
\toprule
Claim & Primary replay or certificate\\
\midrule
\endhead
Quotient census and \(625\)-class dump&
\path{QUOTIENT_CENSUS.md} and \path{census_z37_quotients.cpp};
separate parity audit
\path{audit_z37_quotient_parity.cpp}; all-quotient ambient replay
\path{audit_z37_all_quotient_ambient.py}; packaged
\path{z37_quotient_census_canonical_625.txt},
\path{z37_quotient_census_canonical_625.sha256}, and
\path{verify_package.py}.\\[0.5mm]
Trace law, finite-field counts, first moment, and formal completion&
\path{README.md}, \path{verify_z37_lift_frontier.py}, and
\path{Z37_LIFT_FRONTIER_CERTIFICATE.json}.\\[0.5mm]
Constant rank two&
\path{RANK_TWO_CONJUGATION_OBSTRUCTION.md},
\path{RANK_TWO_JORDAN_OBSTRUCTION.md}, and the corresponding two
verification scripts.\\[0.5mm]
Constant rank three&
\path{rank_three_constant/README.md},
\path{rank_three_constant/audit_rank_three_constant.py}, and
\path{rank_three_constant/RANK_THREE_CONSTANT_CERTIFICATE.json}.\\[0.5mm]
Two physical characteristic-two supports&
\path{char2_support_realization/README.md},
\path{char2_support_realization/verify_all_char2_support.py},
\path{char2_support_realization/TYPE1_SUPPORT_WITNESS.json}, and
\path{char2_support_realization/TYPE2_SUPPORT_WITNESS.json}.\\[0.5mm]
First nonconstant normal form and closures&
\path{first_nonconstant_conjugator/FIRST_NONCONSTANT_CONJUGATOR.md},
\path{first_nonconstant_conjugator/verify_first_nonconstant_gauge.py},
\path{first_nonconstant_conjugator/verify_exceptional_plane_fixed_j.py},
and \path{first_nonconstant_conjugator/FIRST_NONCONSTANT_CERTIFICATE.json}.\\[0.5mm]
Four-cycle and transvection gates&
\path{char2_carry_switches/README.md},
\path{char2_carry_switches/audit_four_cycle_and_carry_invariants.py},
and \path{char2_carry_switches/audit_sparse_unitary_transvections.py}.\\[0.5mm]
Next \(2\)-adic digit&
\path{mod4_support_lift/README.md},
\path{mod4_support_lift/audit_char2_witness_mod4.py}, and
\path{mod4_support_lift/audit_mod4_carry_chart.cpp}.\\
\bottomrule
\end{longtable}
\normalsize

\subsection{Minimal replay}

The package-level certificate check uses only the Python standard library
and is independent of the caller's working directory:
\begin{lstlisting}
python3 hadamard_668_search/papers/semiregular_c37_conference/verify_package.py
\end{lstlisting}

From the repository root:
\begin{lstlisting}
python3 hadamard_668_search/conference_334_z37_lift/verify_z37_lift_frontier.py
python3 hadamard_668_search/conference_334_z37_lift/verify_rank_two_conjugation_obstruction.py
python3 hadamard_668_search/conference_334_z37_lift/verify_rank_two_jordan_obstruction.py
python3 hadamard_668_search/conference_334_z37_lift/char2_support_realization/verify_all_char2_support.py
python3 hadamard_668_search/conference_334_z37_lift/first_nonconstant_conjugator/verify_first_nonconstant_gauge.py
python3 hadamard_668_search/conference_334_z37_lift/first_nonconstant_conjugator/verify_exceptional_plane_fixed_j.py
\end{lstlisting}

The exhaustive quotient replay is:
\begin{lstlisting}
clang++ -O3 -std=c++20 \
  hadamard_668_search/conference_334_z37_lift/census_z37_quotients.cpp \
  -o /tmp/census_z37_quotients

/tmp/census_z37_quotients --dump-canonical \
  > /tmp/z37_625_canonical_final.txt
\end{lstlisting}
Its output must match the packaged dump with SHA-256
\begin{center}
\small
\texttt{c5d8765da49deb39c2ff3407b9d0f265e3ca56c1015d5b0075355c53ca60fb5b}.
\end{center}

The enumerator takes roughly \(45\) seconds and \(60\) MB on an Apple M1
Pro.  The constant-rank-three fixed-\(J\) replay is the longest promoted
single audit, taking about \(5.5\) minutes and \(61\) MB on the same
machine.

The immutable archival location for this package is
\begin{center}
\small
\url{https://github.com/AlecKriebel/Math/releases/tag/h668-research-checkpoint-v1.0.0}.
\end{center}
\section{Conclusion and open boundary}
\label{sec:conclusion}

The semiregular \(C_{37}\) lane survives all quotient-level necessary
conditions and survives characteristic two even after exact margins and
the full \(6/3\) diagonal law are imposed.  This is simultaneously
positive and cautionary.  It produces a sharply defined modular frontier,
but it also shows that characteristic two is not an obstruction.

The stated bounded quotient system has \(625\) permutation classes, two
parity types up to complement, and no zero-diagonal class.  The finite-field
relaxations and the characteristic-\(37\) formal completion explain large
parts of the algebra.  Constant generators through rank three and two
natural first-nonconstant rank-two families cannot realize the binary
diagonal alphabet.  Small local switching also fails on the two frozen
modular supports.

What remains is not a flat search.  The best retained support still has
\(672/1503\) nonzero next-digit carries, and the normalized
characteristic-two relaxation alone is larger than \(2^{720}\).
Any successful continuation needs a global mechanism that couples
physical margins with higher \(2\)-adic carries, or a new obstruction that
acts on the two unitary parity types.

We emphasize the exact endpoint:
\begin{center}
\fbox{\parbox{0.84\textwidth}{\centering
No conference graph or \(H(668)\) is constructed.  Neither the
semiregular \(C_{37}\) lane nor unrestricted existence is excluded.}}
\end{center}

\section*{Computational and AI assistance}

The derivations, programs, certificates, literature audit, and manuscript
were developed with heavy assistance from ChatGPT 5.6 Sol under Alec
Kriebel's direction.  Alec Kriebel is the named human author and describes
himself as a non-specialist who cannot independently certify every
derivation or implementation.  All promoted claims are tied to explicit
local artifacts and exact replay programs listed in
Section~\ref{sec:reproducibility}.  Those artifacts strengthen
auditability; they do not replace expert review, establish priority, or
constitute a proof-assistant formalization.  No external communication was
made in preparing this draft.

{\footnotesize
\begin{thebibliography}{99}
\setlength{\itemsep}{0.25em}

\bibitem{behbahani-lam}
M.~Behbahani and C.~Lam,
\newblock Strongly regular graphs with non-trivial automorphisms,
\newblock \emph{Discrete Mathematics} \textbf{311} (2011), 132--144.
\newblock \url{https://doi.org/10.1016/j.disc.2010.10.005}.

\bibitem{brouwer-srg}
A.~E.~Brouwer,
\newblock Parameters of strongly regular graphs, \(v=301,\ldots,350\),
\newblock online table, accessed 25 July 2026.
\newblock \url{https://aeb.win.tue.nl/graphs/srg/srgtab301-350.html}.

\bibitem{cati-pasechnik}
M.~Cati and D.~V.~Pasechnik,
\newblock A database of constructions of Hadamard matrices,
\newblock arXiv:2411.18897, 2024; revised 2025.
\newblock \url{https://arxiv.org/abs/2411.18897}.

\bibitem{epoch-hadamard}
Epoch AI,
\newblock Hadamard matrix existence,
\newblock FrontierMath open-problem record, accessed 25 July 2026.
\newblock \url{https://epoch.ai/frontiermath/open-problems/hadamard}.

\bibitem{goethals-seidel}
J.~M.~Goethals and J.~J.~Seidel,
\newblock Orthogonal matrices with zero diagonal,
\newblock \emph{Canadian Journal of Mathematics} \textbf{19} (1967),
1001--1010.
\newblock \url{https://doi.org/10.4153/CJM-1967-091-8}.

\bibitem{martinez}
L.~Mart\'inez,
\newblock Strongly regular \(m\)-Cayley circulant graphs and digraphs,
\newblock \emph{Ars Mathematica Contemporanea} \textbf{8} (2015),
195--213.
\newblock
\url{https://www.dlib.si/details/URN:NBN:SI:doc-TCR0ZT1K?language=eng}.

\bibitem{mathon}
R.~Mathon,
\newblock Symmetric conference matrices of order \(pq^2+1\),
\newblock \emph{Canadian Journal of Mathematics} \textbf{30} (1978),
321--331.
\newblock \url{https://doi.org/10.4153/CJM-1978-029-1}.

\bibitem{mckay}
B.~D.~McKay,
\newblock Isomorph-free exhaustive generation,
\newblock \emph{Journal of Algorithms} \textbf{26} (1998), 306--324.
\newblock \url{https://doi.org/10.1006/jagm.1997.0898}.

\bibitem{ramos-hulak-queiroz}
A.~F.~Ramos, D.~B.~Hulak, and R.~J.~G.~B.~de~Queiroz,
\newblock Multiplier obstructions for Legendre pairs of length \(333\),
\newblock arXiv:2607.20765, submitted 22 July 2026.
\newblock \url{https://arxiv.org/abs/2607.20765}.

\bibitem{taylor-classical-groups}
D.~E.~Taylor,
\newblock \emph{The Geometry of the Classical Groups},
\newblock Sigma Series in Pure Mathematics, vol.~9, Heldermann Verlag,
Berlin, 1992.
\newblock \url{https://www.maths.usyd.edu.au/u/don/papers/gcg.pdf}.

\end{thebibliography}
}

\end{document}
